\chapter*{Resumen}
\addcontentsline{toc}{chapter}{Resumen}

Se estudia qué características socioeconómicas, demográficas y sanitarias de los condados de
Estados Unidos se asocian con su tasa de mortalidad por cáncer (muertes anuales por
\num{100000} habitantes), a partir del fichero oficial \var{practica.sav}
(\res{n-condados} condados), que se comprueba celda a celda que coincide con el conjunto
público del que procede \parencite{rippner2017}. Un catálogo de \res{n-reglas} reglas de
validación detecta valores centinela, edades registradas en meses, tamaños de hogar divididos
por cien y una variable que filtra la respuesta; once decisiones documentadas los resuelven y,
tras la depuración, ninguna regla queda en estado de error.

El modelo final se estima por mínimos cuadrados ordinarios (MCO), como pide la orientación,
con \res{n-final} condados y \res{k-final} regresores: \res{n-seleccionadas} términos elegidos
por eliminación hacia atrás con contrastes robustos, \res{n-confusoras} variable reincorporada
por confusión y \res{n-interacciones} interacciones con la región censal. Explica el
\res{pct-r2-final}\,\% de la variabilidad ($R^2 = \res{r2-final}$, $\bar R^2 = \res{r2aj-final}$)
con un error cuadrático medio de \res{rmse-final} muertes por \num{100000} habitantes dentro de
la muestra. Como los errores son heterocedásticos (su varianza decrece con la población del
condado) y dependen dentro de cada estado (correlación intraclase \res{icc-final}), la
inferencia usa errores típicos robustos por conglomerados de estado; los clásicos se presentan
como salida tipo SPSS y una sintaxis generada automáticamente reproduce el modelo en ese
programa.

A igualdad del resto de variables, la mortalidad crece con la incidencia de cáncer y con la
proporción de residentes que sólo tienen cobertura pública, y decrece con el porcentaje de
adultos con título universitario, con un gradiente educativo que es menor en el Noreste. Las
conclusiones se mantienen en las \res{n-especificaciones} especificaciones del análisis de
sensibilidad y en el bootstrap por conglomerados (ponderación por la función
de varianza, efectos fijos de estado, modelo mixto, regresión robusta, exclusión de observaciones
influyentes e imputación múltiple, entre otras). Como complemento, un modelo
predictivo (\res{modelo-elegido}) alcanza un RMSE de \res{test-rmse-elegido} en una partición de
prueba no usada para ajustar ni elegir.

\par\medskip\noindent\textbf{Palabras clave:} regresión lineal múltiple, mínimos cuadrados
ordinarios, errores típicos robustos por conglomerados, diagnóstico de hipótesis, datos
ausentes, mortalidad por cáncer.
